\documentclass[11pt,a4paper]{article}

% ------------------------------------------------------------------
% Paket dasar
% ------------------------------------------------------------------
\usepackage[utf8]{inputenc}
\usepackage[T1]{fontenc}
\usepackage{lmodern}
\usepackage[a4paper,margin=2.5cm,headheight=15pt]{geometry}
\usepackage{amsmath}
\usepackage{booktabs}
\usepackage{tabularx}
\usepackage{longtable}
\usepackage{array}
\usepackage{enumitem}
\usepackage{listings}
\usepackage{xcolor}
\usepackage[most]{tcolorbox}
\usepackage{titlesec}
\usepackage{fancyhdr}
\usepackage[hidelinks]{hyperref}

% ------------------------------------------------------------------
% Nama bagian dalam bahasa Indonesia
% ------------------------------------------------------------------
\renewcommand{\contentsname}{Daftar Isi}
\renewcommand{\tablename}{Tabel}
\renewcommand{\figurename}{Gambar}
\renewcommand{\refname}{Referensi}

% ------------------------------------------------------------------
% Tampilan
% ------------------------------------------------------------------
\definecolor{reportblue}{HTML}{1F4E79}
\definecolor{reportgray}{HTML}{F2F5F8}
\definecolor{codegray}{HTML}{F7F7F7}
\definecolor{codecomment}{HTML}{5A7A5A}

\titleformat{\section}
  {\Large\bfseries\color{reportblue}}
  {\thesection}{0.7em}{}
\titleformat{\subsection}
  {\large\bfseries\color{reportblue}}
  {\thesubsection}{0.6em}{}
\titleformat{\subsubsection}
  {\normalsize\bfseries}
  {\thesubsubsection}{0.5em}{}

\setlength{\parindent}{0pt}
\setlength{\parskip}{0.55em}
\setlength{\emergencystretch}{2em}
\renewcommand{\tabularxcolumn}[1]{>{\raggedright\arraybackslash}m{#1}}

\pagestyle{fancy}
\fancyhf{}
\fancyhead[L]{\small Laporan Mingguan}
\fancyhead[R]{\small Teknik Kompilator}
\fancyfoot[C]{\thepage}
\renewcommand{\headrulewidth}{0.4pt}
\renewcommand{\footrulewidth}{0pt}

\setlist[itemize]{topsep=0.35em,itemsep=0.2em,parsep=0.2em}
\setlist[enumerate]{topsep=0.35em,itemsep=0.25em,parsep=0.2em}

% ------------------------------------------------------------------
% Gaya listing kode
% ------------------------------------------------------------------
\lstdefinestyle{reportcode}{
  backgroundcolor=\color{codegray},
  basicstyle=\ttfamily\small,
  breaklines=true,
  breakatwhitespace=false,
  columns=fullflexible,
  keepspaces=true,
  frame=single,
  rulecolor=\color{black!20},
  framesep=6pt,
  xleftmargin=4pt,
  xrightmargin=4pt,
  showstringspaces=false,
  commentstyle=\color{codecomment}\itshape,
  keywordstyle=\color{reportblue}\bfseries,
  tabsize=2
}
\lstset{style=reportcode}

\newtcolorbox{infobox}[1][]{
  colback=reportgray,
  colframe=reportblue,
  boxrule=0.6pt,
  arc=2pt,
  left=8pt,
  right=8pt,
  top=6pt,
  bottom=6pt,
  #1
}

% ------------------------------------------------------------------
% Judul
% ------------------------------------------------------------------
\title{
  \vspace{-1.5cm}
  {\Large\scshape Teknik Kompilator}\\[0.35cm]
  {\Huge\bfseries Laporan Mingguan}\\[0.35cm]
  {\LARGE \textit{Function} dan \textit{Procedure}}
}
\author{}
\date{25 November 2026}

\begin{document}

% ==================================================================
% Halaman judul
% ==================================================================
\pagenumbering{roman}
\hypersetup{pageanchor=false}
\begin{titlepage}
\centering
\thispagestyle{empty}

\vspace*{2.2cm}
{\Large\scshape Teknik Kompilator\par}

\vspace{1.8cm}
{\Huge\bfseries Laporan Mingguan\par}
\vspace{0.7cm}
{\LARGE \textit{Function} dan \textit{Procedure}\par}

\vspace{2.2cm}
\begin{tabular}{>{\bfseries}rl}
Tanggal     & : 25 November 2026 \\[0.35cm]
\end{tabular}

\vfill
%{\large Laporan ini mencakup kemajuan, hambatan,\\
%dan rencana kerja untuk minggu berikutnya.\par}

\vspace{1.2cm}
\end{titlepage}
\setcounter{page}{2}
\hypersetup{pageanchor=true}

% ==================================================================
% Daftar isi
% ==================================================================
\tableofcontents
\newpage

\pagenumbering{arabic}

% ==================================================================
% 1. Kemajuan
% ==================================================================
\section{Kemajuan}

Target minggu ini adalah menambah pengujian regresi otomatis, memperbaiki penanganan error sintaksis, menguji kombinasi scope yang lebih dalam, serta perapian dan dokumentasi.

\begin{table}[ht]
\centering
\caption{Status target minggu ini}
\label{tab:target}
\begin{tabularx}{\linewidth}{@{}p{0.46\linewidth}p{0.16\linewidth}X@{}}
\toprule
\textbf{Target} & \textbf{Status} & \textbf{Catatan} \\
\midrule
Menambah pengujian regresi otomatis
  & Tidak Selesai & \\
Memperbaiki penanganan error sintaksis
  & Tidak Selesai & \\
Menguji kombinasi scope yang lebih dalam
  & Selesai & \\
Perapian dan dokumentasi
  & Tidak Selesai & \\
\bottomrule
\end{tabularx}
\end{table}

\subsection{Menguji kombinasi scope yang lebih dalam}

Berikut kode untuk menguji nested procedure:


File: \texttt{tes6.txt}
\begin{verbatim}
{
  proc tari(n : integer) {
    proc dindin {
      ;
      put "Dindin badindin"
    }
    ;
    repeat
      dindin
      put "Ooo"
      dindin
      n := n - 1
    until n <= 0
  }
  ;
  tari(4)
}
\end{verbatim}

Berikut output setelah di-compile dan dijalankan di mesin:

\begin{verbatim}
Dindin badindin
Ooo
Dindin badindin
Dindin badindin
Ooo
Dindin badindin
Dindin badindin
Ooo
Dindin badindin
Dindin badindin
Ooo
Dindin badindin
\end{verbatim}

\section{Hambatan}

Untuk minggu ini tidak ada hambatan yang signifikan.

\section{Rencana Minggu Depan}

Belum ada rencana untuk minggu depan.

\appendix
\section{Link Repository}

Berikut link repository GitHub pengerjaan kelompok kami: \url{https://github.com/skrjmcdc/badakdanikan}

\end{document}
